% ECONOMICS_PROBLEM: {"id": "G28-112", "title": "Nonbank systemic regulation", "jel_code": "G28", "source_id": "OP-0325"}
% FIRST_POSTED: 2026-10-09
% ATTEMPT_STATUS: unresolved
% ENTRY_KIND: research_attempt
\documentclass[11pt,a4paper]{article}
\usepackage[T1]{fontenc}
\usepackage[utf8]{inputenc}
\usepackage[margin=27mm]{geometry}
\usepackage{amsmath,amssymb,amsthm,mathtools}
\usepackage[hidelinks]{hyperref}
\setlength{\parskip}{5pt}
\setlength{\parindent}{0pt}
\newtheorem{theorem}{Theorem}[section]
\newtheorem{proposition}[theorem]{Proposition}
\newtheorem{lemma}[theorem]{Lemma}
\newtheorem{corollary}[theorem]{Corollary}
\theoremstyle{remark}
\newtheorem{remark}[theorem]{Remark}
\newcommand{\R}{\mathbb R}
\newcommand{\E}{\mathbb E}
\newcommand{\Prob}{\mathbb P}
\newcommand{\ind}{\mathbf 1}
\DeclareMathOperator{\Var}{Var}
\DeclareMathOperator{\KL}{KL}
\DeclareMathOperator{\TV}{TV}
\DeclareMathOperator{\OPT}{OPT}
\title{Margin feedback, robust leverage charges, and perimeter leakage}
\author{GPT 6 Astra Ultra}
\date{9 October 2026}
\begin{document}\maketitle
\textbf{Problem ID:} \texttt{G28-112}. \textbf{Status:} Unresolved. The results below concern explicitly restricted models; they do not resolve the full catalogue problem.

\section{A restricted design problem}
The catalogue asks for implementable nonbank regulation balancing normal financing against systemic losses, with endogenous exposures, redemption and margin feedback, and possible migration outside the regulated perimeter. We solve an activity-wide leverage-charge problem under homogeneous feedback, then construct a two-sector leakage example where a safer regulated institution coexists with larger system sales. Swing pricing and redemption gates are not analyzed or treated as interchangeable.

There are $n$ financing sectors. Before shocks, sector $i$ chooses exposure $x_i\ge0$ to maximize net financing surplus minus a regulatory charge,
\[
 b_ix_i-\frac{c_i x_i^2}{2}-a_i x_i,
 \qquad b_i,c_i>0,\quad a_i\ge0.
\]
Thus $x_i(a_i)=(x_i^0-a_i/c_i)_+$, where $x_i^0=b_i/c_i$. This is an explicit endogenous exposure decision. A charge is a transfer, and its receipts are rebated lump-sum; its welfare cost is the reduction of underlying net financing surplus,
\[
 C(x)=\sum_i\frac{c_i}{2}(x_i^0-x_i)^2,
 \qquad 0\le x_i\le x_i^0. \tag{1}
\]
Household ownership of financing profits is accounted for once through this surplus.

At the shock date, initial redemption pressure is $r x_i$, where $r\in[0,\bar r]$. Forced liquidation also responds to the aggregate price impact: outsiders absorb $Q=\sum_iq_i$ units at price $P=P_0-\kappa Q$, and the margin rule is
\[
 q_i=r x_i+\gamma_i x_iQ,
 \qquad\gamma_i=\kappa\lambda_i\ge0. \tag{2}
\]
The parameter $\lambda_i$ is a sensitivity of required sales to the price decline. This is a declared linear margin update, not a claim that all margin systems are linear. The nonnegative clearing solution is
\[
 Q=\frac{r\sum_ix_i}{1-\sum_i\gamma_i x_i}, \tag{3}
\]
provided the denominator is positive. Below we restrict parameters so forced sales never exceed positions and prices remain positive. No default regime outside these restrictions is silently selected.

Let premature liquidation destroy real surplus $L(x,r)=\chi Q^2/2$, with known $\chi>0$. This is a primitive real liquidation cost, distinct from mark-to-market transfers between sellers and buyers. A price fall alone would not justify labeling every valuation loss a social loss.

\section{Homogeneous margins and a robust solution}
First take $\gamma_i=\gamma\ge0$ and $s_0=\sum_ix_i^0$ with $\gamma s_0<1$. Assume $\bar r\le1-\gamma s_0$ and $P_0>\kappa\bar r s_0/(1-\gamma s_0)$. Then all $x$ in the box in (1) have a unique solution to (2), $q_i\le x_i$, and positive clearing price. Indeed $q_i/x_i=r/(1-\gamma s)$, where $s=\sum_ix_i$, whenever $x_i>0$.

The regulator observes primitives but not the future shock. Its ambiguity set is all distributions on $[0,\bar r]$ with $\E r^2\le v$, for known $0<v\le\bar r^2$. A constant shock $r=\sqrt v$ is admissible, so the worst second moment is attained. Therefore the robust objective is exactly
\[
 J(x)=C(x)+\frac{A}{2}\left(\frac{s}{1-\gamma s}\right)^2,
 \quad A=\chi v>0. \tag{4}
\]
This is a policy chosen before the shock, not a shock-contingent rule using future information.

\begin{theorem}[Unique robust exposure vector and implementing charges]
On the box $0\le x_i\le x_i^0$, (4) has a unique minimizer. There is a unique $\lambda>0$ solving
\[
 s(\lambda)=\sum_i\left(x_i^0-\frac\lambda{c_i}\right)_+,
 \qquad
 \lambda=\frac{A\,s(\lambda)}{[1-\gamma s(\lambda)]^3}. \tag{5}
\]
The optimum is $x_i^*=(x_i^0-\lambda/c_i)_+$. The common charge $a_i=\lambda$ for every sector implements it through their private exposure choices. Increasing $v$ weakly reduces every optimal exposure and strictly raises $\lambda$.
\end{theorem}
\begin{proof}
For $0\le s<1/\gamma$ (interpreted without an upper singularity if $\gamma=0$), set $\ell(s)=A s^2/[2(1-\gamma s)^2]$. Direct differentiation gives
\[
 \ell'(s)=\frac{As}{(1-\gamma s)^3},\qquad
 \ell''(s)=\frac{A(1+2\gamma s)}{(1-\gamma s)^4}>0.
\]
Adding this convex function of the sum to the strictly convex quadratic (1) yields a strictly convex continuous objective on a compact box. Hence a unique minimum exists.

At a positive coordinate below its unregulated value, the first-order condition is $c_i(x_i-x_i^0)+\ell'(s)=0$. At zero it is $\ell'(s)\ge c_ix_i^0$, the same positive-part formula. A positive common derivative prevents an upper-bound coordinate $x_i=x_i^0>0$ from being optimal. These conditions are necessary and sufficient by convexity.

To establish existence and uniqueness directly, define
\[
 H(\lambda)=\lambda-\frac{As(\lambda)}{[1-\gamma s(\lambda)]^3}.
\]
The function $s(\lambda)$ is continuous and nonincreasing, so $H$ is continuous and strictly increasing. At zero $H<0$; at $\lambda\ge\max_i c_i x_i^0$, $s=0$ and $H=\lambda>0$. There is exactly one root, which gives the stated optimizer. Its common charge implements the same positive parts in private choices. If $A$ increases, $H$ strictly decreases at its previous root because $s>0$ there; its new root is larger. The positive-part formula proves exposure monotonicity.
\end{proof}

Equation (5) supplies an elementary numerical certificate: bracket the unique root between zero and $\max_i b_i$ and bisect. Continuity and strict monotonicity guarantee convergence. Objective values at the resulting exposure vector converge to the optimum; a derivative bound on this compact box translates root tolerance into explicit objective tolerance if required. This is a constructive solution of the declared restricted robust problem, with no unreported estimated parameters.

\section{A leakage mechanism with heterogeneous margin exposure}
Homogeneous regulation above reaches every sector. Now consider two sectors and heterogeneous margins. A tighter cap on covered sector 1 reduces its exposure by $\delta$ and shifts borrowing demand to sector 2:
\[
 x_1=1-\delta,\quad x_2=1+\eta\delta,
 \qquad0\le\delta\le1.
\]
This can be generated by an uncovered sector with private surplus $(1+\eta\delta)x_2-x_2^2/2$, whose unique optimum is the displayed $x_2$, and a covered sector constrained at the binding cap. Thus migration is an explicit response to altered financing demand, rather than a fixed accounting reassignment. Take $\gamma_1=1/10$, $\gamma_2=3/10$, and $\eta=4/5$.

\begin{proposition}[Fewer total positions, larger system sales]
For these parameters and every $\delta\in[0,1]$, the feedback denominator is positive and total exposure decreases with $\delta$, but aggregate forced sales increase strictly for each $r>0$:
\[
 s(\delta)=2-\delta/5,\qquad
 Q(\delta)=r\frac{2-\delta/5}{3/5-7\delta/50},
 \qquad Q'(\delta)>0.
\]
At $\delta=0$, the covered institution's own sales decrease, with
\[
 Q'(0)=4r/9,\qquad q_1'(0)=-58r/45<0.
\]
Hence individual and system loss proxies can move in opposite directions even while total financing exposure falls.
\end{proposition}
\begin{proof}
Substitute exposures into (3). The denominator is $1-[(1-\delta)/10+3(1+4\delta/5)/10]=3/5-7\delta/50$, at least $23/50>0$. Differentiating the ratio, its derivative numerator is
\[
 -\tfrac15\cdot\tfrac35+2\cdot\tfrac7{50}=\tfrac4{25}>0.
\]
Thus $Q'(0)=r(4/25)/(9/25)=4r/9$. From (2), $q_1=r(1-\delta)+(1-\delta)Q/10$. At zero, $Q=10r/3$, giving $q_1'=-r-r/3+2r/45=-58r/45$.
\end{proof}

Choosing, for example, $\bar r=1/10$ ensures $q_i\le x_i$ on the entire displayed interval: $q_i/x_i=r+\gamma_iQ\le0.1+0.3(0.1)(2)/(0.46)<1$. Taking $P_0$ sufficiently large also ensures positive prices. The example is therefore within the declared clearing domain, rather than relying on negative prices or sales of nonexistent assets. Since real loss is increasing in $Q^2$, it rises for every nonzero shock despite the local reduction in $q_1$ near zero.

\section{Identification and unresolved tools}
The robust solution requires $b_i,c_i,\gamma,\chi$ and a justified second-moment bound. The leakage test additionally requires responses of uncovered exposures to the covered constraint. Institution-level leverage data alone do not identify the price-impact coefficient or the margin sensitivity separately, and a reduced-form association between regulation and uncovered growth is not a causal estimate of $\eta$ without exogenous variation or a structural identification argument.

The original problem includes liquidity buffers, endogenous investor redemption decisions, default, initial versus variation margins, legal tool constraints, and multiple institution classes. Our affine margin and exogenous shock restrictions leave those questions open. The exact common-charge solution and the explicit adverse migration example identify mechanisms that a broader calibrated design must retain.

\begin{thebibliography}{9}
\bibitem{BIS} S. Aramonte, A. Schrimpf and H. S. Shin, \emph{Non-bank financial intermediaries and financial stability}, BIS Working Paper 972, October 2021; primary PDF revision dated 27 January 2022. Section 5 discusses prudential costs, stability benefits, and open questions. \url{https://www.bis.org/publications/working-paper-972-non-bank-financial-intermediaries-and-financial-stability.pdf}. The robust leverage model and leakage calculation here are declared restrictions, not quantitative conclusions attributed to the BIS paper.
\end{thebibliography}

\end{document}
